\appendix
\section{Appendix}
\subsection{Robust Risk Bounds and the Certified Barycentric Family}
\label{app:barycentric_bound}

This appendix provides the complete derivation underlying
Proposition~\ref{prop:barycentric_family}. The argument adapts the
CF-consistent Wasserstein robust-bound construction to the present
post-hoc refinement setting.

\begin{proof}[Proof of Proposition~\ref{prop:barycentric_family}]
Recall the uncertainty lens
\begin{equation}
\mathcal{C}
=
\left\{
\mu:
W_2(\mu,\mathbb{P}_1)\le \epsilon,\;
W_2(\mu,\mathbb{P}_{\mathrm{CF}})\le \delta
\right\},
\end{equation}
and the corresponding worst-case risk
\begin{equation}
R(\mathbb{Q})
=
\sup_{\mu\in\mathcal{C}}
W_2^2(\mathbb{Q},\mu).
\end{equation}

Fix any $\mu\in\mathcal{C}$. By the triangle inequality,
\begin{align}
W_2(\mathbb{Q},\mu)
&\le
W_2(\mathbb{Q},\mathbb{P}_{\mathrm{CF}})
+
W_2(\mathbb{P}_{\mathrm{CF}},\mu)
\nonumber\\
&\le
W_2(\mathbb{Q},\mathbb{P}_{\mathrm{CF}})+\delta,
\label{eq:app_cf_triangle}
\end{align}
and similarly,
\begin{align}
W_2(\mathbb{Q},\mu)
&\le
W_2(\mathbb{Q},\mathbb{P}_1)
+
W_2(\mathbb{P}_1,\mu)
\nonumber\\
&\le
W_2(\mathbb{Q},\mathbb{P}_1)+\epsilon.
\label{eq:app_target_triangle}
\end{align}
Therefore,
\begin{equation}
W_2^2(\mathbb{Q},\mu)
\le
\min\left\{
\left(W_2(\mathbb{Q},\mathbb{P}_{\mathrm{CF}})+\delta\right)^2,\,
\left(W_2(\mathbb{Q},\mathbb{P}_1)+\epsilon\right)^2
\right\}.
\label{eq:app_min_bound}
\end{equation}

For any $\alpha\in[0,1]$ and any $a,b\ge 0$,
\begin{equation}
\min\{a,b\}
\le
(1-\alpha)a+\alpha b.
\end{equation}
Applying this inequality to Eq.~\ref{eq:app_min_bound} gives
\begin{align}
W_2^2(\mathbb{Q},\mu)
\le\;&
(1-\alpha)
\left(
W_2(\mathbb{Q},\mathbb{P}_{\mathrm{CF}})+\delta
\right)^2
\nonumber\\
&+
\alpha
\left(
W_2(\mathbb{Q},\mathbb{P}_1)+\epsilon
\right)^2.
\label{eq:app_convex_bound}
\end{align}

We next control the cross terms. For any $\gamma>0$ and $a,b\ge 0$,
Young's inequality gives
\begin{equation}
2ab
\le
\gamma a^2+\frac{b^2}{\gamma}.
\label{eq:young}
\end{equation}
Hence, for arbitrary $\gamma_0,\gamma_1>0$,
\begin{align}
\left(
W_2(\mathbb{Q},\mathbb{P}_{\mathrm{CF}})+\delta
\right)^2
\le\;&
(1+\gamma_0)
W_2^2(\mathbb{Q},\mathbb{P}_{\mathrm{CF}})
\nonumber\\
&+
\left(1+\frac{1}{\gamma_0}\right)\delta^2,
\label{eq:app_cf_young}
\end{align}
and
\begin{align}
\left(
W_2(\mathbb{Q},\mathbb{P}_1)+\epsilon
\right)^2
\le\;&
(1+\gamma_1)
W_2^2(\mathbb{Q},\mathbb{P}_1)
\nonumber\\
&+
\left(1+\frac{1}{\gamma_1}\right)\epsilon^2.
\label{eq:app_target_young}
\end{align}

Combining Eqs.~\ref{eq:app_convex_bound}--\ref{eq:app_target_young},
we obtain, for every $\mu\in\mathcal{C}$,
\begin{align}
W_2^2(\mathbb{Q},\mu)
\le\;&
(1-\alpha)(1+\gamma_0)
W_2^2(\mathbb{Q},\mathbb{P}_{\mathrm{CF}})
\nonumber\\
&+
\alpha(1+\gamma_1)
W_2^2(\mathbb{Q},\mathbb{P}_1)
+
K(\alpha,\gamma_0,\gamma_1),
\label{eq:app_upper_bound}
\end{align}
where
\begin{align}
K(\alpha,\gamma_0,\gamma_1)
:=\;&
(1-\alpha)
\left(1+\frac{1}{\gamma_0}\right)\delta^2
\nonumber\\
&+
\alpha
\left(1+\frac{1}{\gamma_1}\right)\epsilon^2.
\label{eq:app_constant}
\end{align}
Importantly, $K(\alpha,\gamma_0,\gamma_1)$ is independent of
$\mathbb{Q}$.

Since Eq.~\ref{eq:app_upper_bound} holds for every
$\mu\in\mathcal{C}$, taking the supremum over the lens yields
\begin{align}
R(\mathbb{Q})
\le\;&
(1-\alpha)(1+\gamma_0)
W_2^2(\mathbb{Q},\mathbb{P}_{\mathrm{CF}})
\nonumber\\
&+
\alpha(1+\gamma_1)
W_2^2(\mathbb{Q},\mathbb{P}_1)
+
K(\alpha,\gamma_0,\gamma_1).
\label{eq:app_robust_bound}
\end{align}

Define
\begin{equation}
A
:=
(1-\alpha)(1+\gamma_0),
\qquad
B
:=
\alpha(1+\gamma_1).
\label{eq:app_AB}
\end{equation}
For $A+B>0$, the $\mathbb{Q}$-dependent part of
Eq.~\ref{eq:app_robust_bound} can be written as
\begin{align}
&A\,W_2^2(\mathbb{Q},\mathbb{P}_{\mathrm{CF}})
+
B\,W_2^2(\mathbb{Q},\mathbb{P}_1)
\nonumber\\
&\qquad =
(A+B)
\left[
(1-\lambda)
W_2^2(\mathbb{Q},\mathbb{P}_{\mathrm{CF}})
+
\lambda
W_2^2(\mathbb{Q},\mathbb{P}_1)
\right],
\label{eq:app_normalization}
\end{align}
where
\begin{equation}
\lambda
:=
\frac{B}{A+B}
=
\frac{\alpha(1+\gamma_1)}
{(1-\alpha)(1+\gamma_0)+\alpha(1+\gamma_1)}.
\label{eq:app_lambda}
\end{equation}

Because the additive term $K$ is independent of $\mathbb{Q}$ and
the multiplicative factor $A+B$ is positive, minimizing the robust
upper bound with respect to $\mathbb{Q}$ is equivalent to minimizing
\begin{equation}
(1-\lambda)
W_2^2(\mathbb{Q},\mathbb{P}_{\mathrm{CF}})
+
\lambda
W_2^2(\mathbb{Q},\mathbb{P}_1).
\label{eq:app_barycentric_objective}
\end{equation}
Thus, for every admissible choice of the bounding parameters, the
corresponding robust upper bound is minimized by a $2$-Wasserstein
barycenter of $\mathbb{P}_{\mathrm{CF}}$ and $\mathbb{P}_1$ with
weights $(1-\lambda,\lambda)$.

It remains to show that the induced barycentric weight spans the
entire interval $[0,1]$. Fix arbitrary $\gamma_0,\gamma_1>0$ and
define
\begin{equation}
\lambda(\alpha)
=
\frac{\alpha(1+\gamma_1)}
{(1-\alpha)(1+\gamma_0)+\alpha(1+\gamma_1)}.
\end{equation}
This mapping is continuous on $[0,1]$, with
\begin{equation}
\lambda(0)=0,
\qquad
\lambda(1)=1.
\end{equation}
Moreover,
\begin{equation}
\frac{d\lambda(\alpha)}{d\alpha}
=
\frac{(1+\gamma_0)(1+\gamma_1)}
{
\left[
(1-\alpha)(1+\gamma_0)
+
\alpha(1+\gamma_1)
\right]^2
}
>0.
\label{eq:app_lambda_derivative}
\end{equation}
Hence $\lambda(\alpha)$ is strictly increasing and maps $[0,1]$
onto $[0,1]$.

Consequently, the minimizers of this family of valid robust upper
bounds span the Wasserstein barycentric family
\begin{equation}
\mathcal{Q}_{\mathrm{bar}}
=
\left\{
\mathbb{Q}(\lambda):
\lambda\in[0,1]
\right\},
\end{equation}
which proves Proposition~\ref{prop:barycentric_family}.
\end{proof}

\subsection{Proof of the Monotone Transport Geometry}
\label{app:transport_geometry}

\begin{proof}[Proof of Proposition~\ref{prop:transport_geometry}]
Let
\[
t=\eta\lambda^\star.
\]
By Eq.~\ref{eq:partial_ot_refinement},
$\mathbb{P}_{\mathrm{rCF}}^{(\eta)}$ is the displacement interpolation
at time $t$ induced by an optimal coupling between
$\mathbb{P}_{\mathrm{CF}}$ and $\mathbb{P}_1$.

Since displacement interpolation under the quadratic Wasserstein
metric defines a constant-speed geodesic, for any $t\in[0,1]$,
\begin{align}
W_2\!\left(
\mathbb{P}_{\mathrm{CF}},\mathbb{Q}(t)
\right)
&=
t\,W_2(\mathbb{P}_{\mathrm{CF}},\mathbb{P}_1),\\
W_2\!\left(
\mathbb{Q}(t),\mathbb{P}_1
\right)
&=
(1-t)\,W_2(\mathbb{P}_{\mathrm{CF}},\mathbb{P}_1).
\end{align}
Substituting $t=\eta\lambda^\star$ and
$D=W_2(\mathbb{P}_{\mathrm{CF}},\mathbb{P}_1)$ gives
\begin{align}
W_2\!\left(
\mathbb{P}_{\mathrm{rCF}}^{(\eta)},
\mathbb{P}_{\mathrm{CF}}
\right)
&=
\eta\lambda^\star D,\\
W_2\!\left(
\mathbb{P}_{\mathrm{rCF}}^{(\eta)},
\mathbb{P}_1
\right)
&=
(1-\eta\lambda^\star)D.
\end{align}
The first quantity is nondecreasing in $\eta$, while the second is
nonincreasing in $\eta$. Summing the two identities yields
Eq.~\ref{eq:geometric_frontier}.
\end{proof}


\subsection{Proof of Proposition~\ref{prop:future_validity}}
\label{app:future_validity}

\begin{proof}
Since $f$ is $L_f$-Lipschitz, the surrogate
$\ell_f(z)=[\tau-f(z)]_+$ is also $L_f$-Lipschitz. Hence, by the
Kantorovich--Rubinstein inequality,
\begin{equation}
\mathbb{E}_{z\sim\mathbb{Q}}[\ell_f(z)]
\le
\mathbb{E}_{z\sim\mu_f}[\ell_f(z)]
+
L_f W_1(\mathbb{Q},\mu_f).
\end{equation}
Using $W_1\le W_2$ and
Assumption~\ref{ass:future_valid_reference},
\begin{equation}
\mathbb{E}_{z\sim\mathbb{Q}}[\ell_f(z)]
\le
\alpha_f
+
L_f W_2(\mathbb{Q},\mu_f).
\end{equation}
Since $\mu_f\in\mathcal{C}$ and
\[
R(\mathbb{Q})
=
\sup_{\mu\in\mathcal{C}}
W_2^2(\mathbb{Q},\mu),
\]
we have
\[
W_2(\mathbb{Q},\mu_f)
\le
\sqrt{R(\mathbb{Q})}.
\]
Therefore,
\begin{equation}
\mathbb{E}_{z\sim\mathbb{Q}}[\ell_f(z)]
\le
\alpha_f
+
L_f\sqrt{R(\mathbb{Q})},
\end{equation}
which proves the result.
\end{proof}


